\section{Feedback strictly improves a two-flip minimax subproblem}
\label{sec:feedback-counterexample}
\noindent\textbf{PROVED-ON-PAPER with certified finite rational arithmetic.}
This is a counterexample to removing cross-coin acquisition feedback without
loss. It is not a separation under expected budgets, not an asymptotic rate
separation, and not a lower bound for the full class $\mathcal Q_\rho$.
No novelty or formalization claim is made.

Fix $\rho=1/4$ and the known magnitude law
$\mu=\tfrac12\delta_{1/4}+\tfrac12\delta_{1/8}$.
The magnitude of each actual coin remains hidden. Consider the admissible
two-point parameter class
$\mathcal H=\{q_0\equiv0,q_1\equiv1/4\}$, with targets
$\theta_0=0$ and $\theta_1=1/4$.
All genuine coins are iid and every requested flip is charged. The
following comparison imposes a \emph{deterministic total cap of two flips}.

The noninteractive comparison class preselects its independent local coin
policies using an arbitrary shared seed, independent of all coin tapes.
Each local policy may stop using its own observations and the seed, but its
acquisition choices cannot depend on any other coin's observations. The
terminal estimator may be an arbitrary nonlinear function of all acquired
data and the seed. The sum of the preselected policies' actual flips must
obey the cap almost surely at every parameter; an external budget monitor
that cancels one policy because of another's observations is cross-coin
feedback and is outside this comparison class.

\begin{theorem}\label{feedback:strict}
For this two-point problem, the minimax squared error of the noninteractive
comparison class is strictly greater than $15313/10^6$.
There is a feedback procedure with exactly two flips and maximum squared
error strictly less than $7623/500000$.
In particular its minimax squared-error advantage exceeds $67/10^6$.
\end{theorem}

\paragraph{Likelihoods and the two fixed schedules.}
For a word $h$ of length $n$ with $k$ heads, its one-coin probability under
hypothesis $j$ is the rational number
\[
 f_j(h)=\frac12\sum_{d\in\{1/4,1/8\}}
 \left[(1-j/4)(1/2-d)^k(1/2+d)^{n-k}
       +(j/4)(1/2+d)^k(1/2-d)^{n-k}\right].
\]
If two bits $(x,y)$ come from the same coin, their probability is
$p_{j,O}(x,y)=f_j(xy)$. For two fresh coins it is
$p_{j,N}(x,y)=f_j(x)f_j(y)$.
These probabilities use only actual observed bits, not conditional gap
information supplied to the estimator.

Every member of the stated noninteractive class is statistically dominated
by a random mixture of these two full schedules. To see this, condition
on all independent seeds, including local randomization. If a local policy
can request two flips with positive probability under $q_0$, no other
policy can request even one with positive probability under $q_0$;
otherwise independence would make total work greater than two with positive
probability. If no policy can request two, at most two policies can request
one. Finite words have strictly positive probability under both hypotheses,
so the same statement holds for both hypotheses and for every potentially
queried finite history. Unused pre-created identities reveal no data and
can be omitted. A one-coin policy that stops early is dominated by completing
two flips of that coin; policies with at most one flip each are dominated
by observing one flip from each of two fresh coins. The original estimator
can ignore the additional bits. This also handles zero observations.
The chosen dominating schedule depends only on the seed, not on the
hypothesis. Randomization therefore cannot lower its Bayes risk below
the smaller Bayes risk of the two schedules.

\paragraph{A feedback policy and its exact estimator.}
Flip one fresh coin. After a tail, use the second flip on a fresh coin.
After a head, revisit the original coin. In symbols, its probability of
$(x,y)$ is
\[
 p_{j,A}(x,y)=
 \begin{cases}f_j(x)f_j(y),&x=0,\\ f_j(xy),&x=1.\end{cases}
\]
This costs exactly two flips in both hypotheses. Fix the rational number
$\pi=1075683009/2147483648$, interpreted only as a design constant, and
output the posterior mean
\[
 \widehat\theta_A(x,y)
 =\frac14\frac{\pi p_{1,A}(x,y)}
 {(1-\pi)p_{0,A}(x,y)+\pi p_{1,A}(x,y)}.
\]
The estimator is fully specified and takes values in $(0,1/4)$; there is
no uncharged pilot, latent-gap query, or unknown-parameter calibration.

\paragraph{Finite rational verification.}
For any schedule $P\in\{O,N,A\}$ the optimal Bayes squared error under
prior $\pi$ is exactly
\[
 \mathcal B_P(\pi)=\frac1{16}\pi(1-\pi)
 \sum_{x,y\in\{0,1\}}
 \frac{p_{0,P}(x,y)p_{1,P}(x,y)}
 {(1-\pi)p_{0,P}(x,y)+\pi p_{1,P}(x,y)}.
\]
Let $r_{j,A}=\sum_{x,y}p_{j,A}(x,y)
(\widehat\theta_A(x,y)-j/4)^2$.
The script \texttt{feedback\_certificate.py} computes all these fractions
exactly, checks that each likelihood vector sums to one, and checks the
independent posterior-variance identity for every Bayes risk. It verifies
\[
 \max_j r_{j,A}<\frac{7623}{500000}
 <\frac{15313}{10^6}<\min\{\mathcal B_O(\pi),\mathcal B_N(\pi)\}.
\]
For orientation only, the adaptive maximum is approximately
$0.01524525010158912$, while the smaller fixed-schedule Bayes lower bound
is approximately $0.015313875487111883$.
The inequalities, rather than these decimal displays, are the certificate.
Every minimax risk is at least its Bayes risk under the specified prior,
so the displayed chain proves Theorem~\ref{feedback:strict}.

The script also gives exact certificates for magnitude laws with equal
weights at $(1/4,1/64)$ and $(1/4,1/1024)$, with separately pinned dyadic
design priors. Thus strict feedback gain is not specific to neighboring
gap scales. It remains a finite cap-two effect: none of the certificates
establishes a growing advantage, a critical-law lower bound, or the same
comparison for algorithms having only a uniform expected-work budget.

\paragraph{Why this matters, and the remaining gap.}
An unrestricted estimator cannot in general be replaced at no loss by a
preselected independent acquisition design and an optimal nonlinear
postprocessor. The first bit changes whether information should be acquired
from the same latent coin or a fresh one, even in this tiny admissible
subproblem. The certificate does not imply that feedback defeats the
critical logarithmic obstruction in the worst case over arbitrary $q$.
Amplifying the finite advantage into such a statement requires a hard
family or construction whose advantage grows with the number of gap scales,
while controlling the original uniform expected budget. No such
amplification is proved here.
